% ============================================================
% 6_Conclusion.tex
% ============================================================

\section{Conclusion}
\label{sec:conclusion}

We presented a comprehensive distribution-driven experimental study of child
speech emotion recognition, spanning 63 controlled comparisons across three
datasets, three pooling strategies, and seven experiment series. Rather than
claiming a single universally superior architecture, our results establish
\textit{when} and \textit{why} specific design choices matter for children's
SER.

\textbf{This work demonstrates} that the statistical distribution gap between
children's and adults' emotional speech---measurable via Fr\'{e}chet Distance
in WavLM feature space---is the dominant factor governing SER accuracy, more
so than architectural details. The FD--WA monotonic relationship holds across
age, style, and enhancement dimensions, providing a quantitative diagnostic for
anticipating deployment performance.

\textbf{The decisive evidence} comes from four converging findings:
(1)~Prosody-Guided pooling shows child-specific efficacy (SA$-$PG gap: 2.0pp
on children vs.\ 10.0pp on adults); (2)~augmenting children's data with adult
speech causes catastrophic degradation ($-$25pp); (3)~zero-shot cross-corpus
transfer uniformly fails ($<$34\% WA); and (4)~learnable layer fusion provides
the largest single-module gain (+10pp over single-layer baselines) with
consistent mid-to-upper layer preference across all corpora.

\textbf{The broader implication} is that child SER cannot be treated as
``adult SER with higher pitch.'' It requires corpus design, augmentation
strategies, and architectural priors that are explicitly matched to the target
pediatric population. Distribution awareness must be elevated from a
post-hoc diagnostic to a first-class design constraint.

\textbf{Boundary conditions:} Our conclusions are drawn from three datasets
with limited linguistic and cultural diversity (Malay, German, English).
Extension to tonal languages, non-WEIRD child populations, and clinical
pediatric speech (e.g., autism, speech disorders) requires further validation.
The FD diagnostic, while effective within our WavLM-based framework, may
produce different absolute values under alternative SSL backbones.

We release all experimental configurations, model checkpoints, and evaluation
code to facilitate reproducible child SER research and encourage the community
to adopt distribution-aware evaluation protocols.
